 \paragraph*{04.05.09.02 Zur Evaluation und Auswahl von Lieferanten und Partnern beitragen}
\kcilabel{para:04.05.09.02_Lieferantenauswahl}{04.05.09.02}
\label{para:04.05.09.02_Lieferantenauswahl}


% Task
\par Meine Aufgabe bestand darin, die Evaluation und Auswahl von Lieferanten und Partnern für die Beschaffung externer, nicht durch den \gls{gu}-Vertrag gedeckter Leistungen im Rahmen von \gls{r2h} aktiv zu unterstützen. Dabei war es entscheidend, geeignete technische und organisatorische Auswahlkriterien zu definieren, potenzielle Lieferanten zu identifizieren und vergleichbar zu bewerten sowie in Zusammenarbeit mit der \gls{tpl} und dem \gls{zek} eine belastbare Entscheidungsgrundlage zu schaffen.

% Action
\par Um die Evaluation und Auswahl von Lieferanten und Partnern zu unterstützen, entwickelte ich zunächst ein Bewertungsraster mit gewichteten Kriterien (technische Fähigkeiten in \gls{s4}, Projekterfahrung, Verfügbarkeit, Referenzen und vertragliche Passfähigkeit). Anschließend analysierte ich die im \gls{adac} vorhandenen Lieferanten mit \gls{rv}, um geeignete Kandidaten zu identifizieren. Die beiden identifizierten Lieferanten kontaktierte ich anschließend und bat sie um die \gls{cv} geeigneter Kandidaten. In enger Zusammenarbeit mit der \gls{tpl} und dem \gls{zek} wurden die \gls{cv} gesichtet, vergleichend bewertet und in eine fundierte Entscheidungsgrundlage für die Lieferantenauswahl überführt.

% Result
\par Die Evaluation und Auswahl der Lieferanten führte letztlich zur Einbindung von drei Programmierern aus zwei unterschiedlichen externen Dienstleistern mit umfassendem \gls{s4}-Know-how, die die Implementierung der Welle 1.1 ff. im Rahmen von \gls{r2h} unterstützten. Durch die strukturierte Auswahl konnten die erforderliche Expertise sichergestellt und externe Leistungen termingerecht bereitgestellt werden. Dadurch wurden Umsetzungsrisiken im kritischen Pfad reduziert und ein wesentlicher Beitrag zum Projekterfolg geleistet.

% Reflect
\par Die sorgfältige Evaluation und Auswahl der Lieferanten hat gezeigt, dass eine strukturierte, kriterienbasierte Vorgehensweise entscheidend für den Projekterfolg ist. Insbesondere die frühe Einbindung der relevanten Stakeholder (\gls{tpl}, \gls{zek}) und die transparente Vergleichbarkeit der Anbieter verbesserten die Entscheidungsqualität. Durch die gezielte Auswahl externer Partner konnte die erforderliche Expertise gesichert und die termingerechte Umsetzung der Projektziele nachhaltig unterstützt werden.